\chapter{Solver Architecture and the Per-Step Pipeline}\label{ch:arch}

This chapter provides the map to which the remainder of the manual refers. We
describe how the continuous problem of \cref{eq:lagrangian-ode} is reduced to a
discrete pipeline, the principal data structures that the pipeline operates on,
the one-time assembly that precedes the time loop, and the ordered sequence of
operations performed at each step. Later chapters expand individual stages of this
pipeline; here we are concerned only with their arrangement.

\section{From the continuous equation to a discrete pipeline}\label{sec:arch-pipeline}

The engine advances \cref{eq:lagrangian-ode} by an explicit time-stepping scheme:
time is divided into steps of length $\dt$, and at each step the velocity is
sampled along the path and used to displace the particle. Evaluating the velocity
at the particle position decomposes into a fixed sequence of operations that is
the same at every step and for every scheme. The particle must first be
\emph{located} on the unstructured mesh --- the nearest node and the containing
element must be identified --- because the interpolation stencils are expressed
in terms of mesh entities. The velocity is then \emph{interpolated in space}, from
the surrounding nodes or element centres to the particle position, and
\emph{reconstructed in time}, from the bracketing stored snapshots to the substep
time required by the integrator. The integrator \emph{advances} the position, and
finally the \emph{boundary conditions} are applied: a particle that has left the
domain is flagged and its step undone, while a particle that has reached the
surface or the sea floor is clamped. The result is recorded at a prescribed
cadence.

This decomposition --- locate, interpolate in space, reconstruct in time,
integrate, apply boundary conditions, record --- is realized directly in the
source and is depicted in \cref{fig:pipeline}. The chapters of
\crefrange{ch:mesh}{ch:lifecycle} correspond one-to-one to its stages.

\section{Principal data structures}\label{sec:arch-data}

Three objects carry the state of a run. The \fn{Mesh} holds the static geometry:
the node coordinates, the element connectivity, the per-node sea-floor depth, the
vertical interfaces, the frame and radius, and the search structures and adjacency tables built once and
cached on it. The \fn{VelocityField} holds the model fields for a single
time level, already shaped according to where they live (at nodes, at element
centres, or at interfaces) and tagged with their layout. The \fn{ParticleState}
holds the evolving per-particle quantities: position, status, the cached nearest
node, the most recently evaluated velocity, and --- when the run uses them --- the failure reason, the
release step and the carried element. \Cref{tab:datastructures}
summarizes their principal members; the chapters indicated develop each in detail.

\begin{table}[t]
\centering
\caption{The three principal data structures and their principal members. Shapes
are given for $N$ nodes, $E$ elements, $K$ vertical interfaces (hence $K-1$
layers), and $P$ particles.}
\label{tab:datastructures}
\small
\begin{tabularx}{\linewidth}{@{}l l X@{}}
\toprule
Structure & Member & Meaning \\
\midrule
\fn{Mesh}
  & \code{nodes\_lon/lat} & node coordinates, degrees ($N$); see \cref{ch:mesh} \\
  & \code{elem\_nodes}    & element connectivity $(E,3)$ or $(E,4)$, $-1$-padded \\
  & \code{elem\_center\_lon/lat} & element centroids ($E$) \\
  & \code{depth}          & per-node sea-floor depth, positive down ($N$) \\
  & \code{sigma}         & vertical interfaces ($K$): depths on z-level meshes, $\sigma$ values on sigma meshes \\
  & \code{vertical\_coord} & \code{zlevel} or \code{sigma} \\
  & \code{node\_to\_elems} & node$\rightarrow$element adjacency; see \cref{ch:localization} \\
  & \code{rotation\_angles} & frame of the coordinates, or \code{None}; see \cref{ch:frame} \\
  & (cached)              & grid hash and node table, element neighbours, areas, coastal flags \\
\midrule
\fn{VelocityField}
  & \code{u}, \code{v}    & horizontal velocity; node $(N)$/$(K{-}1,N)$ or element $(E)$/$(K{-}1,E)$ \\
  & \code{w}              & vertical velocity at interfaces $(K,N)$, or \code{None} (2-D) \\
  & \code{eta}            & sea-surface height ($N$) \\
  & \code{uv\_layout}     & \code{"node"} or \code{"element"}; selects the stencil \\
  & \code{scalars}        & sampled fields (e.g.\ temperature, salinity) \\
  & \code{kh\_*}, \code{az} & horizontal and vertical diffusivities, when used \\
\midrule
\fn{ParticleState}
  & \code{lon}, \code{lat}, \code{z} & position ($P$) \\
  & \code{status}         & lifecycle state; see \cref{ch:lifecycle} \\
  & \code{nearest\_node}  & cached localization result ($P$) \\
  & \code{u}, \code{v}, \code{w} & last evaluated velocity ($P$) \\
  & \code{failure\_reason}, \code{release\_step}, \code{elem} & optional: cause of death, release step, carried element \\
\bottomrule
\end{tabularx}
\end{table}

\codeanchor{The data structures are defined in \srcfile{core/mesh.py} (\fn{Mesh}),
\srcfile{core/io.py} (\fn{VelocityField}) and \srcfile{core/particles.py}
(\fn{ParticleState}).}

\section{Run assembly}\label{sec:arch-assembly}

Before the time loop begins, a single assembly phase prepares the run, in this order. The reader
appropriate to the mesh type is selected and the mesh is loaded (\cref{ch:mesh}); the exact-anchor guard is
set and the grid-hash cell is chosen --- pinned, measured or the reader's (\cref{sec:loc-cell}) --- and in three
dimensions the vertical interfaces are trimmed to \conf{max\_levels}. The forcing is opened lazily, so that
snapshots are read from disk only as they are needed, and the release positions are read. The time axis is
resolved: the interval between stored snapshots is detected, the step is set, and for time-averaged forcing
each timestamp is shifted from the end of its averaging window to its midpoint (\cref{ch:input}). Only then is
the release schedule built, since release steps are offsets divided by the step; it pre-allocates the full
particle array and freezes not-yet-released particles (\cref{ch:lifecycle}). The seeds are rotated into the
computational frame (\cref{ch:frame}). A set of probes then checks, before any output file exists, the
combinations that would fail later --- the sampled fields, the velocity layout against the interpolation scheme,
the diffusion, beaching, periodic and adaptive-scheme settings --- so that a misconfigured run fails in
seconds; a few remain for the first step (spatially varying diffusion in three dimensions, the adaptive
scheme on a pixel mesh with the exact guard off). For each realization the output writer is opened with its provenance attributes (\cref{ch:output}),
the shared snapshot cache is attached when configured (\cref{sec:input-shared}), and the time loop runs.

\codeanchor{The assembly order is implemented in \fn{run\_tracking}
(\srcfile{tracking.py}); the time loop in \fn{run\_single\_realization}
(\srcfile{tracking.py}).}

\section{The per-step loop}\label{sec:arch-loop}

\begin{figure}[t]
\centering
\begin{tikzpicture}[node distance=3.6mm and 16mm,flowbox/.append style={text width=54mm},sidebox/.append style={text width=54mm}]
  \node[flowbox] (act)   {Activate due releases\\[1pt]\code{activate\_released}};
  \node[flowbox,below=of act]   (vel)   {Build velocity payload\\[1pt]\footnotesize bracket $\rightarrow$ stencil $\rightarrow$ cache $\rightarrow$ lazy view};
  \node[flowbox,below=of vel]   (int)   {Integrate one step: \code{step\_fn}\\[1pt]\footnotesize Euler\,/\,RK2\,/\,RK4\,/\,RK45};
  \node[flowbox,below=of int]   (fold)  {Periodic fold (if set)};
  \node[flowbox,below=of fold]  (beach) {Beaching hook (if enabled)};
  \node[flowbox,below=of beach] (clock) {Advance clock\\[1pt]\footnotesize interpolated mode: $t\leftarrow t+\tdir\,\dt$};
  \node[flowbox,below=of clock] (rec)   {Record \& checkpoint\\[1pt]\footnotesize global step $g$: $(g{+}1)$ multiple of the cadence};

  \draw[flowarrow] (act)--(vel);
  \draw[flowarrow] (vel)--(int);
  \draw[flowarrow] (int)--(fold);
  \draw[flowarrow] (fold)--(beach);
  \draw[flowarrow] (beach)--(clock);
  \draw[flowarrow] (clock)--(rec);
  \draw[looparrow] (rec.east) -- ++(0.9,0) |- node[pos=0.25,sloped,below,font=\footnotesize,text=lagblue!60!black]{next step} (act.east);

  \node[sidebox,right=22mm of int] (sub) {%
    \textbf{Inside \code{step\_fn}, per particle:}\\[2pt]
    1.\ exact nearest node at each site\\
    2.\ \emph{(3-D)} find vertical layer\\
    3.\ interpolate $u,v$ (with temporal blend)\\
    4.\ \emph{(3-D)} interpolate $\Wvel$\\
    5.\ advance position (every stage)\\
    6.\ re-locate; test containment\\
    7.\ commit, or revert, flag \code{OUTSIDE}\\
    \phantom{7.\ }and record the failure reason\\
    8.\ \emph{(3-D)} clamp $\depx$ to surface/floor\\
    9.\ diffusive kicks, re-checked};
  \draw[-{Stealth[length=4pt]},gray] (sub.west) -- (int.east);
\end{tikzpicture}
\caption{The per-step pipeline executed by \fn{run\_single\_realization}. The main loop (left) activates any
particles whose release time has arrived, assembles the velocity for the step, integrates, folds particles
back into periodic windows, applies the optional beaching hook, advances the clock, and records and
checkpoints at the configured cadence. The side panel lists the operations performed for each particle
inside the integration step; stages marked \emph{(3-D)} are present only in three-dimensional runs.}
\label{fig:pipeline}
\end{figure}

At each step the engine performs the sequence shown in \cref{fig:pipeline}. It first activates any particles
whose scheduled release time has arrived, so that they are advected during the same step
(\cref{ch:lifecycle}). It then assembles the velocity payload for the step: the stored snapshots that bracket
the current time are identified, the reconstruction stencil is resolved (two snapshots for linear
reconstruction, four for monotone piecewise-cubic), the snapshots are taken from the cache or read, and the
step receives either a snapshot or a lazy view that blends a field only when it is read, together with the
target snapshot for the Runge--Kutta stages (\cref{sec:temp-linear,sec:input-payload}). On a sigma mesh the
surface elevation at the start and the end of the step is supplied as well (\cref{sec:vert-sigma-step}). The
integration step is invoked once; the operations it performs for each particle are listed in the side panel
of \cref{fig:pipeline} and are the subject of \crefrange{ch:localization}{ch:integration}. Particles that left
a periodic window are folded back, and if beaching is enabled a post-step hook evaluates the stranding and
resuspension criteria (\cref{ch:lifecycle}). In the interpolated mode the clock advances by $\tdir\,\dt$; in the
discrete mode the snapshot index advances by $\tdir$, one record forward or backward. The state is recorded, and checkpointed, after every step
whose global index $g$ satisfies that $g+1$ is a multiple of \conf{output\_freq}, respectively
\conf{checkpoint\_freq}, so that a resumed run records the same steps as an uninterrupted one
(\cref{ch:output}); one more checkpoint follows the last executed step. A record is labelled with the time at
the end of its step, samples the requested scalars at that time, un-rotates positions and velocities to
geographic coordinates, writes the rows of not-yet-released particles as missing values, and writes the
velocity as missing wherever the particle was not alive through the whole step.

The boundary conditions are applied inside the step, immediately after the position is advanced. The
containing-element test serves as the domain check: a particle whose new position is not contained by any
element --- whether it has crossed an open boundary or left a coastline --- has its step undone and its status
set to \code{OUTSIDE}. Vertical limits, by contrast, are enforced by clamping rather than removal: a particle
that would cross the surface or the local sea floor is held at that limit (\cref{ch:vertical,ch:lifecycle}).
With the option \code{-v} the loop prints, per step, the time spent reading (\code{io=}), stepping
(\code{step=}), beaching (\code{beach=}) and recording (\code{rec=}) (\cref{sec:perf-measure}).

\codeanchor{The step sequence is the loop body of \fn{run\_single\_realization}
(\srcfile{tracking.py}); the per-particle operations are implemented in the
integration kernels of \srcfile{numerics/integrators.py}.}

\section{Forward and backward integration}\label{sec:arch-direction}

The engine integrates equally in either time direction, and it does so without a
separate backward code path. The direction factor $\tdir\in\{+1,-1\}$ is folded
into the interpolated velocity: every interpolation routine returns
$\tdir\,u$, $\tdir\,v$ and $\tdir\,\Wvel$ rather than the raw model velocity, while
the time step $\dt$ remains positive. Backward integration is therefore a sign
reversal of the sampled velocity, and the same integrator advances the position in
both cases. The step counter increments in both directions, and the simulation
clock advances by $\tdir\,\dt$. With stochastic diffusion the backward run is the adjoint process: the
advection is reversed, the diffusive drift and noise are not (\cref{sec:diff-model}). This symmetry is what the reversibility tests of
\cref{ch:validation} exploit: integrating forward and then backward over the same
interval, with the same scheme, returns a particle to its starting point to within
a small tolerance.

\section{Fused and unfused execution paths}\label{sec:arch-fused}

The per-particle operations of \cref{fig:pipeline} are realized by two interchangeable implementations. A set
of \emph{fused} kernels, compiled just-in-time, performs the whole step --- localization, interpolation,
integration and the boundary check --- within a single compiled loop over the particles, for the common
configurations to which they apply: two-dimensional least-squares runs, time-dependent two-dimensional
barycentric runs on node velocities (RK2 and RK4, linear interpolation in time, no element tracking), three-dimensional Euler with least-squares velocities on z-level meshes (vertical combine \conf{nearest}
or \conf{leastsq}), and
three-dimensional RK2 with node-layout barycentric velocities when the vertical velocity is combined by
\conf{nearest} or \conf{barycentric}. An \emph{unfused} path performs the same operations as separate parallel
passes over all particles and covers every other configuration. Where both exist they produce bit-for-bit
identical trajectories, statuses and failure reasons under the default exact anchor guard
(\conf{mesh.exact\_nearest}); the details are in \cref{ch:fused}.
